% Part: first-order-logic
% Chapter: tableaux
% Section: quantifier-rules

\documentclass[../../../include/open-logic-section]{subfiles}

\begin{document}

\olfileid[hi]{fol}{tab}{qrl}

\olsection{परिमाणक-नियम}

\subsection{$\lforall$ के नियम}

\begin{defish}
\AxiomC{\sFmla{\True}{\lforall[x][!A(x)]}}
\RightLabel{\TRule{\True}{\forall}}
\UnaryInfC{\sFmla{\True}{!A(t)}}
\DisplayProof
\hfill
\AxiomC{\sFmla{\False}{\lforall[x][!A(x)]}}
\RightLabel{\TRule{\False}{\lforall}}
\UnaryInfC{\sFmla{\False}{!A(a)}}
\DisplayProof
\end{defish}

\TRule{\True}{\lforall} में $t$ बंद पद है (अर्थात् ऐसा पद जिसमें कोई चर नहीं
है)। \TRule{\False}{\lforall} में $a$~ऐसा !!a{constant} है जो
\TRule{\False}{\lforall} नियम के ऊपर की शाखा में कहीं भी न आए। $a$ को
\TRule{\False}{\forall} अनुमान का \emph{अभिलाक्षणिक चर} कहते हैं।\footnote{ऊपर
के नियम में $a$ !!a{constant} है, फिर भी ऐतिहासिक कारणों से ``अभिलाक्षणिक
चर'' शब्द प्रयुक्त होता है।}

\subsection{$\lexists$ के नियम}

\begin{defish}
\AxiomC{\sFmla{\True}{\lexists[x][!A(x)]}}
\RightLabel{\TRule{\True}{\lexists}}
\UnaryInfC{\sFmla{\True}{!A(a)}}
\DisplayProof
\hfill
\AxiomC{\sFmla{\False}{\lexists[x][!A(x)]}}
\RightLabel{\TRule{\False}{\lexists}}
\UnaryInfC{\sFmla{\False}{!A(t)}}
\DisplayProof
\end{defish}

यहाँ भी $t$~बंद पद है, और $a$~ऐसा !!a{constant} है जो
\TRule{\True}{\lexists} नियम के ऊपर की शाखा में नहीं आता। $a$ को
\TRule{\True}{\lexists} अनुमान का \emph{अभिलाक्षणिक चर} कहते हैं।

\TRule{\False}{\lforall} या \TRule{\True}{\lexists} अनुमान के ऊपर की शाखा
में अभिलाक्षणिक चर के न आने की शर्त को \emph{अभिलाक्षणिक-चर शर्त} कहते हैं।

\begin{explain}
स्मरण करें: जब $!A$ ऐसा !!a{formula} हो जिसमें !!{variable}~$x$ मुक्त है, तो
इसे~$!A(x)$ लिखकर दिखाते हैं। उसी संदर्भ में $!A(t)$,
संक्षेप में~$\Subst{!A}{t}{x}$ है। अतः \TRule{\False}{\lexists} नियम को ऐसे
भी लिख सकते हैं:
\begin{prooftree}
  \AxiomC{\sFmla{\False}{\lexists[x][!A]}}
  \RightLabel{\TRule{\False}{\lexists}}
  \UnaryInfC{\sFmla{\False}{\Subst{!A}{t}{x}}}
\end{prooftree}
ध्यान दें कि $t$,~$!A$ में पहले से आ सकता है; जैसे $!A$ का रूप
$\Atom{\Obj P}{t,x}$ हो सकता है। इसलिए $\sFmla{\False}{\Atom{\Obj
P}{t,t}}$ को~$ \sFmla{\False}{\lexists[x][\Atom{\Obj P}{t,x}]}$ से निष्कर्षित
करना~\TRule{\False}{\lexists} का सही प्रयोग है। किन्तु
\TRule{\False}{\lforall} और~\TRule{\True}{\lexists} की अभिलाक्षणिक-चर शर्तें
माँगती हैं कि !!{constant}~$a$,~$!A$ में न आए। अतः
$\sFmla{\False}{\Atom{\Obj P}{a,a}}$ को
$\sFmla{\False}{\lforall[x][\Atom{\Obj P}{a,x}]}$ से
~$\TRule{\False}{\lforall}$ द्वारा सही ढंग से निष्कर्षित नहीं कर सकते।
\end{explain}

\begin{explain}
\TRule{\True}{\lforall} और \TRule{\False}{\lexists} में पद~$t$ पर कोई
प्रतिबंध नहीं है। दूसरी ओर, \TRule{\True}{\lexists} और
\TRule{\False}{\lforall} नियमों में अभिलाक्षणिक-चर शर्त माँगती है कि
!!{constant}~$a$ संबंधित अनुमान के ऊपर वाली किसी शाखा में कहीं भी न आए। तंत्र
की सुदृढ़ता सुनिश्चित करने के लिए यह आवश्यक है। इस शर्त के बिना,
$\lexists[x][\formula{A}(x)] \lif \lforall[x][\formula{A}(x)]$ के लिए
निम्नलिखित बंद !!{tableau} होता:
\begin{center}
\begin{tableau}{}
  [\sFmla{\False}{\lexists[x][\formula{A}(x)] \lif \lforall[x][\formula{A}(x)]}, just=\TAss
    [\sFmla{\True}{\lexists[x][\formula{A}(x)]},
      just={\TRule{\False}{\lif}[1]}
      [\sFmla{\False}{\lforall[x][\formula{A}(x)]},
        just={\TRule{\False}{\lif}[1]}
        [\sFmla{\True}{\formula{A}(a)},
          just={\TRule{\True}{\lexists}[2]}
          [\sFmla{\False}{\formula{A}(a)},
            just={\TRule{\False}{\lforall}[3]}, close]
        ]
      ]
    ]
  ]
\end{tableau}
\end{center}
परंतु $\lexists[x][\formula{A}(x)] \lif
\lforall[x][\formula{A}(x)]$ वैध नहीं है।
\end{explain}

\end{document}
